% REV02: proofs consolidated in 02_dualidad_t.tex and 03_pantallas_coordinadas.tex.
\subsection{Exceptional incidence and coordination of the screens}
\label{exc:pantallas-dualidad}

The direction selected by \(K\) and the Leech lattice arise from distinct
representations of the constructed incidence. The former determines the
orthogonal reduction \(12\to11\to10\); the latter enters the isotropic
reduction \(26\to24\). The coordination preserves the preceding space, its
relations, and the information transmitted by each projection.

The sectoral representation on T-duality and dimensional screens jointly
develops this coordination. Its first block brings together the involution of
integer coordinates, unitary conjugation with its domains, and the change of
residual section with quotient memory. The second constructs the morphism
\(\mathcal R_M^{\mathrm{alg}}\) and proves the isomorphism of the quotient
presentations. The graph preserves the coordinate in \(V_{11}\) together
with its projection onto \(V_{10}\), while the isotropic quotient recovers
the Leech lattice. The respective kernels and codomains are thus preserved.

The action ellipse is constructed before duality is applied and fixes its
dimensionless radius together with the reciprocal radius. The compact sector
then links exceptional incidence to a determined quadratic form. The fields,
action, and operators of the subsequent physical realization are presented
in their proper domain, preserving this algebraic antecedent and its internal
proofs.
